\subsection{معماری و پیاده‌سازی}
کلاس اصلی \lr{LoLAAttention} در ماژول \lr{Codes/framework/attentions/lola\_attention.py} قرار داشته و اصلاحات مهندسی حساسی به شرح زیر در آن پیاده شده است:
\begin{itemize}
    \item \textbf{مدیریت حافظه در انکودرهای دوطرفه (\lr{LL-FIX-01}):} کش پویا تنها در پردازش‌های دیکودر خودبازگشتی جاری فعال شده و برای مدل‌های انکودر (\lr{DistilBERT} و \lr{RoBERTa}) به طور خودکار غیرفعال می‌گردد تا از آلودگی تانسوری میان دسته‌ها و رشد بی‌رویه مصرف حافظه جلوگیری به عمل آید (کاهش حدود ۷۰ درصدی مصرف حافظه).
    \item \textbf{بازنشانی صریح حافظه کش (\lr{LL-FIX-03}):} متد \lr{cache.reset()} در ابتدای توالی‌های نامرتبط فراخوانی می‌شود تا تانسورهای معلق در حافظه بازیافت گردند.
    \item \textbf{سازگاری کامل با \lr{GQA} و \lr{RoPE}:} پیاده‌سازی توابع انطباق ابعاد سرها و افزودن موقعیت زاویه‌ای با در نظر گرفتن جابه‌جایی تانسورهای فشرده‌شده.
\end{itemize}

\subsection{محدودیت‌ها و تحلیل شکست}
با وجود کارایی بالا در سرعت و حافظه، برخی رفتارهای انحرافی شایان توجه است:
\begin{itemize}
    \item \textbf{اشباع دقت در مدل‌های با رتبه پایه (\lr{PASS\_LOW}):} دقت ۵۰.۹۲٪ در مدل‌های \lr{RoBERTa} و \lr{GPT-2} نشان می‌دهد که تقریب اولیه با رتبه ۳۲ بر روی نگاشت‌های عمیق بدون بازآموزی دقیق لایه‌های بالاتر، سبب از دست رفتن بخش قابل توجهی از اطلاعات رتبه‌های منفرد فرعی ماتریس می‌گردد.
    \item \textbf{تحلیل معماری بر روی \lr{Qwen}:} در مدل \lr{Qwen}، ثبت دقت ۵۲.۰۶٪ حاکی از تثبیت وضعیت در حاشیه پایه است. دلیل این امر، تقابل میان فشرده‌سازی سه‌بخشی کش و نحوه انتشار بردارهای چرخش موقعیتی (\lr{RoPE}) در فرکانس‌های بالا است؛ زیرفضای تصویرشده ارتوگونال در بافر توکن‌های باستانی، تعامد ذاتی فرکانس‌های زاویه‌ای توکن‌های دوردست را تا حدی مغشوش می‌سازد که منجر به تضعیف استدلال‌های حساس به ترتیب کلمات در متون طولانی می‌گردد.
\end{itemize}
